% !TEX root = ../tex/5-5.tex
% 第5章：本文は青、補足は緑
\newcommand{\asmain}{%
    \setbeamercolor{frametitle}{fg=white,bg=main}%
    \setbeamercolor{structure}{fg=main}%
    \setbeamercolor{itemize item}{fg=main}%
    \setbeamercolor{itemize subitem}{fg=main}%
    \setbeamercolor{enumerate item}{fg=main}%
}
\newcommand{\assuppl}{%
    \setbeamercolor{frametitle}{fg=white,bg=examplecol}%
    \setbeamercolor{structure}{fg=examplecol}%
    \setbeamercolor{itemize item}{fg=examplecol}%
    \setbeamercolor{itemize subitem}{fg=examplecol}%
    \setbeamercolor{enumerate item}{fg=examplecol}%
}
% フレームタイトルの「補足」だけを□で囲む
\newcommand{\supplmark}{%
    \setlength{\fboxsep}{1.5pt}%
    \setlength{\fboxrule}{0.7pt}%
    \fbox{\textbf{補足}}%
}
\asmain


%=============================================
\begin{frame}{5.5 Closing Remarks}
    \begin{block}{この章で閉じた第三の柱（論文 5.5）}
        変分・スコアに続く、決定的フローに根ざす拡散モデルの見方。
    \end{block}
    \begin{itemize}
        \item 出発：Normalizing Flows（変数変換による可逆写像）
        \item 連続化：Neural ODE（速度場が変換を支配）
        \item 到達：Flow Matching（シミュレーション不要の速度回帰）
    \end{itemize}
\end{frame}
%=============================================
\asmain
\begin{frame}{第5章の道筋（全体像）}
    \begin{enumerate}
        \item \textbf{5.1 NF / Neural ODE}\\
        離散の可逆変換から、連続時間の速度場へ
        \item \textbf{5.2 Flow Matching}\\
        確率経路と連続の式。条件付き速度への回帰
        \item \textbf{5.3 経路と速度の構成}\\
        ガウス特殊例、path-first / flow-first
        \item \textbf{5.4（任意）}\\
        Canonical affine・Rectify / Reflow
    \end{enumerate}
\end{frame}
%=============================================
\asmain
\begin{frame}{旅程：NF $\to$ NODE $\to$ FM}
    \begin{block}{Normalizing Flows}
        変数変換公式で、単純事前とデータ間の正確な可逆写像を学ぶ。
    \end{block}
    \begin{block}{Neural ODEs}
        変換を連続時間ダイナミクスにする。\\
        代償：訓練ループ内で ODE シミュレーションが必要。
    \end{block}
    \begin{block}{Flow Matching}
        確率経路 $\{p_t\}$ と連続の式を満たす速度を先に定め、\\
        条件付けトリックで周辺速度マッチを条件付き回帰へ落とす。\\
        \textbf{シミュレーション不要}の学習。
    \end{block}
\end{frame}
%=============================================
\asmain
\begin{frame}{条件付けトリックという共通パターン}
    変分・スコアと同様、FM も「条件付け」で学習を手なずける。
    \begin{itemize}
        \item 周辺速度 $v_t$ の直接マッチは一般に困難
        \item 条件付き速度 $u_t(\cdot\mid \cdot)$ は閉じた形で書けることが多い
        \item 条件付き回帰の平均が、周辺目標を誘導する
    \end{itemize}
    \begin{exampleblock}{拡散モデルの位置づけ}
        ガウス事前をデータへ運ぶ決定的フロー学習の、特別な場合として読める。
    \end{exampleblock}
\end{frame}
%=============================================
\asmain
\begin{frame}{Score SDE との対応}
    同じ拡散確率経路に対して：
    \begin{columns}[T,onlytextwidth]
        \begin{column}{0.48\textwidth}
            \begin{block}{Score SDE（第4章）}
                スコア $s_\phi\approx\nabla\log p_t$ を学ぶ。\\
                そこから PF-ODE の速度を組む。
            \end{block}
        \end{column}
        \begin{column}{0.48\textwidth}
            \begin{block}{Flow Matching（第5章）}
                同じ（または一般の）経路の\\
                速度場 $v_\theta$ を直接当てる。
            \end{block}
        \end{column}
    \end{columns}
    \vfill
    \begin{block}{共通点}
        どちらも、扱いにくい周辺目標を\\
        扱いやすい条件付き目標へ落とせる。
    \end{block}
\end{frame}
%=============================================
\asmain
\begin{frame}{三本の柱が揃った}
    \centering
    \renewcommand{\arraystretch}{1.35}
    \begin{tabular}{lll}
        \textbf{視点} & \textbf{起源} & \textbf{学ぶ量の例} \\
        \hline
        変分 & VAE / DDPM & ノイズ予測・ELBO \\
        スコア & EBM / NCSN & $\nabla\log p_t$ \\
        フロー & NF / NODE / FM & 速度場 $v_t$ \\
    \end{tabular}
    \vfill
    \begin{itemize}
        \item 出自は違うが、いずれも連続時間の生成過程へ収束
        \item 学習を可能にする条件付け戦略が繰り返し現れる
    \end{itemize}
\end{frame}
%=============================================
\asmain
\begin{frame}{章頭の到達目標との対応}
    章頭の ``After this chapter, you will be able to...'' への対応。
    \begin{itemize}
        \item NF / NODE / FM を「確率の輸送」でつなぐ
        \item Score SDE：スコア $\to$ PF-ODE。FM：速度を直接
        \item 連続の式が、ODE速度と密度発展をつなぐ
        \item ガウス拡散を超える一般ソース／ターゲット輸送
        \item 同じ周辺・異なる速度／カップリング／粒子軌道
        \item path-first と flow-first の二通りの条件付き構成
    \end{itemize}
\end{frame}
%=============================================
\assuppl
\begin{frame}{\supplmark：到達目標チェックリスト}
    \footnotesize
    \begin{enumerate}
        \item 離散可逆変換 $\to$ 連続速度場として輸送を語れるか
        \item 同じ経路で「スコア経由」と「速度直接」を対比できるか
        \item 連続の式を、密度レベルでの拘束として説明できるか
        \item FMが一般分布間でも simulation-free な理由を言えるか
        \item $\{p_t\}$ と、それを実現する速度・カップリング・軌道を区別できるか
        \item 条件付き経路から速度、または条件付き流マップから速度、を説明できるか
    \end{enumerate}
\end{frame}
%=============================================
\asmain
\begin{frame}{核心の一文}
    \begin{block}{Flow Matchingの核心}
        確率経路を先に定め、連続の式を満たす速度を\\
        条件付き回帰で当てる。生成は決定的ODEの数値解である。
    \end{block}
    \begin{itemize}
        \item NODEの訓練内シミュレーションを避ける
        \item 拡散は、ガウス事前$\to$データのフロー学習の一例
        \item Score SDEのPF-ODEと同じ舞台を、別の学習量で見る
    \end{itemize}
\end{frame}
%=============================================
\asmain
\begin{frame}{第6章へ：統一レンズ}
    次章では三視点を、ひとつの枠にまとめる。
    \begin{enumerate}
        \item 変分・スコア・フローが、根本で数学的に同等であること
        \item Fokker--Planckが、三視点に共通する密度発展の法則であること
    \end{enumerate}
    \begin{block}{共通レシピの予告}
        前向きの破壊過程で周辺経路を定め、\\
        その経路を運ぶ時間依存ベクトル場を学ぶ。
    \end{block}
\end{frame}
%=============================================
\asmain
\begin{frame}{第5章のまとめ}
    \begin{itemize}
        \item NF $\to$ NODE $\to$ FM で、決定的確率輸送の第三の柱を見た
        \item FMは条件付き速度回帰により、simulation-free に学習する
        \item Score SDEと同じ確率流を、速度を直接当てる見方で読む
        \item （任意）Rectify / Reflow で周辺軌道を真っ直ぐに近づける
        \item 次は第6章：三視点の統一と Fokker--Planck
    \end{itemize}
\end{frame}
